\lecture{8}{2026-10-06}{More Discrete Logarithm Algorithms}

In this lecture, we explore two non-generic algorithms for computing discrete logarithms: the \emph{Pohlig--Hellman algorithm}, which exploits the factorization of the group order, and the \emph{index calculus algorithm}, which leverages the factor base structure of \(\mathbb{Z}_p^\times\) to achieve subexponential running time.

\section*{The Pohlig--Hellman Algorithm}

Let \(G\) be a cyclic group of order \(n\) generated by \(g\), and let \(h \in G\).
The goal is to find the discrete logarithm \(x \in \mathbb{Z}_n\) such that
\[
    g^x = h.
\]
The Pohlig--Hellman algorithm~\cite{PohligHellman1978} reduces the problem of finding \(x\) modulo \(n\) to finding \(x\) modulo the prime-power factors of \(n\), and then combining the results via the Chinese Remainder Theorem.

\subsection*{Prime-Power Order Groups}

Suppose first that the group order is a prime power, \(n = p^e\), where \(p\) is prime and \(e \geq 1\).
We represent the unknown exponent \(x \in \{0, 1, \dots, p^e - 1\}\) in base \(p\):
\[
    x = x_0 + x_1 p + x_2 p^2 + \dots + x_{e-1} p^{e-1}, \quad \text{where } x_i \in \{0, 1, \dots, p-1\}.
\]
The digits \(x_0, x_1, \dots, x_{e-1}\) are determined sequentially:

\begin{enumerate}
    \item \textbf{Base step for \(x_0\):}
    Raise both sides of \(g^x = h\) to the power \(p^{e-1}\):
    \begin{align*}
        h^{p^{e-1}} &= (g^x)^{p^{e-1}} \\
        &= g^{(x_0 + x_1 p + \dots + x_{e-1} p^{e-1}) p^{e-1}} \\
        &= (g^{p^{e-1}})^{x_0} \cdot (g^{p^e})^{x_1 + x_2 p + \dots} \\
        &= (g^{p^{e-1}})^{x_0},
    \end{align*}
    since \(g^{p^e} = 1\).
    Let \(\alpha = g^{p^{e-1}}\) and \(\beta_0 = h^{p^{e-1}}\).
    Then \(\alpha\) has prime order \(p\), and \(\beta_0 = \alpha^{x_0}\).
    We can compute \(x_0 = \log_\alpha(\beta_0)\) using any DLP algorithm in a group of order \(p\) (such as baby-step giant-step or Pollard's rho) in \(O(\sqrt{p})\) operations.

    \item \textbf{Inductive step for \(x_k\):}
    Assume inductively that \(x_0, x_1, \dots, x_{k-1}\) have been determined.
    Define
    \[
        y_k = \sum_{j=0}^{k-1} x_j p^j.
    \]
    Then \(x - y_k = x_k p^k + x_{k+1} p^{k+1} + \dots + x_{e-1} p^{e-1}\).
    Dividing \(h\) by \(g^{y_k}\) and raising to the power \(p^{e-1-k}\), we obtain
    \[
        \beta_k = (h \cdot g^{-y_k})^{p^{e-1-k}} = (g^{x - y_k})^{p^{e-1-k}} = (g^{p^{e-1}})^{x_k} \cdot (g^{p^e})^{x_{k+1} + \dots} = \alpha^{x_k}.
    \]
    Again, this is a discrete logarithm problem in the subgroup of order \(p\) generated by \(\alpha\), which yields \(x_k \in \{0, \dots, p-1\}\) in \(O(\sqrt{p})\) group operations.
\end{enumerate}

After \(e\) iterations, the full exponent \(x \equiv \sum_{k=0}^{e-1} x_k p^k \pmod{p^e}\) is recovered.

\begin{algorithm}[H]
    \caption{Pohlig--Hellman for Prime-Power Order}\label{alg:pohlig-hellman-prime-power}
    \begin{algorithmic}[1]
        \Input generator \(g \in G\) of order \(p^e\), element \(h \in \langle g \rangle\), prime \(p\), exponent \(e \geq 1\)
        \Output integer \(x \in \{0, \dots, p^e - 1\}\) such that \(g^x = h\)
        \State set \(\alpha \gets g^{p^{e-1}}\)
        \State set \(x \gets 0\)
        \For{\(k = 0\) \textbf{to} \(e - 1\)}
            \State set \(\beta \gets (h \cdot g^{-x})^{p^{e-1-k}}\)
            \State solve \(\alpha^{x_k} = \beta\) for \(x_k \in \{0, \dots, p-1\}\) \Comment{via Pollard's rho or BSGS}
            \State \(x \gets x + x_k p^k\)
        \EndFor
        \State \Return \(x\)
    \end{algorithmic}
\end{algorithm}

\begin{remark}[Complexity for Prime-Power Order]
    The algorithm performs \(e\) discrete logarithm computations in a subgroup of prime order \(p\).
    Using Pollard's rho algorithm, each step takes \(O(\sqrt{p})\) group operations, resulting in an overall time complexity of \(O(e \sqrt{p})\) group operations.
\end{remark}

\subsection*{General Pohlig--Hellman Algorithm}

Now consider the general case where the group order \(n\) factors into distinct prime powers as
\[
    n = p_1^{e_1} p_2^{e_2} \cdots p_k^{e_k}.
\]
For each \(i \in \{1, \dots, k\}\), define
\[
    n_i = \frac{n}{p_i^{e_i}}, \quad g_i = g^{n_i}, \quad h_i = h^{n_i}.
\]
Notice that \(g_i\) has order \(p_i^{e_i}\), and
\[
    h_i = (g^x)^{n_i} = (g^{n_i})^x = g_i^x = g_i^{x \pmod{p_i^{e_i}}}.
\]
Thus, we can apply \Cref{alg:pohlig-hellman-prime-power} with base \(g_i\), target \(h_i\), and group order \(p_i^{e_i}\) to find
\[
    a_i \equiv x \pmod{p_i^{e_i}}.
\]
Having computed \(a_i\) for all \(i \in \{1, \dots, k\}\), we obtain a system of congruences:
\[
    x \equiv a_i \pmod{p_i^{e_i}} \quad \text{for } i = 1, \dots, k.
\]
Because the moduli \(p_1^{e_1}, \dots, p_k^{e_k}\) are pairwise coprime, the Chinese Remainder Theorem guarantees a unique solution for \(x \pmod{n}\).

\begin{algorithm}[H]
    \caption{General Pohlig--Hellman Algorithm}\label{alg:pohlig-hellman-general}
    \begin{algorithmic}[1]
        \Input generator \(g \in G\) of order \(n = p_1^{e_1} \cdots p_k^{e_k}\), element \(h \in \langle g \rangle\)
        \Output unique integer \(x \in \{0, \dots, n - 1\}\) such that \(g^x = h\)
        \For{\(i = 1\) \textbf{to} \(k\)}
            \State set \(g_i \gets g^{n / p_i^{e_i}}\), \(h_i \gets h^{n / p_i^{e_i}}\)
            \State \(a_i \gets \textsc{Pohlig--Hellman-Prime-Power}(g_i, h_i, p_i, e_i)\)
        \EndFor
        \State use the Chinese Remainder Theorem to find \(x \in \{0, \dots, n-1\}\) such that \(x \equiv a_i \pmod{p_i^{e_i}}\) for all \(i\)
        \State \Return \(x\)
    \end{algorithmic}
\end{algorithm}

\begin{remark}[Cryptographic Implications]
    The overall complexity of the general Pohlig--Hellman algorithm is
    \[
        O\left( \sum_{i=1}^k e_i (\log n + \sqrt{p_i}) \right) \text{ group operations.}
    \]
    Consequently, computing discrete logarithms in \(G\) is only as hard as computing discrete logarithms in its largest prime-order subgroup.
    If \(n\) is smooth (i.e., all prime factors \(p_i\) are small), DLP is easy.
    To resist the Pohlig--Hellman attack, cryptographic cyclic groups must be chosen such that their order is divisible by a large prime factor (or has prime order).
\end{remark}

\section*{Index Calculus}

Unlike generic algorithms that treat the group as a black box, the \emph{index calculus algorithm} is a non-generic method that exploits the concrete representation of elements in \(\mathbb{Z}_p^\times\) and the unique factorization of integers.
It runs in subexponential time.

\subsection*{Smooth Numbers and Factor Base}

The core idea is that integers with small prime factors are easy to factor.

\begin{definition}[$B$-Smooth Number]
    An integer \(m \geq 1\) is called \emph{\(B\)-smooth} if all prime factors of \(m\) are less than or equal to \(B\).
\end{definition}

Fix a smoothness bound \(B\).
The set of all primes up to \(B\) forms the \emph{factor base}:
\[
    \mathcal{F} = \{p_1, p_2, \dots, p_k\},
\]
where \(p_1 = 2, p_2 = 3, p_3 = 5, \dots\), and \(p_k \leq B\) is the largest prime not exceeding \(B\).
By the Prime Number Theorem, \(k = \pi(B) \approx B / \ln B\).

\subsection*{Algorithm Description}

Let \(g\) be a primitive root modulo \(p\), and let \(h \in \mathbb{Z}_p^\times\).
We wish to compute \(x = \log_g(h) \pmod{p-1}\).
The index calculus algorithm proceeds in two main stages:

\paragraph{Stage 1: Relation Collection and Factor Base Discrete Logarithms}
We search for exponents \(x_j \in \mathbb{Z}_{p-1}\) such that the least positive residue of \(g^{x_j} \pmod{p}\) is \(B\)-smooth.
When a smooth residue is found, it factors completely over the factor base:
\begin{align*}
    g^{x_1} &\equiv p_1^{e_{11}} p_2^{e_{12}} \cdots p_k^{e_{1k}} \pmod{p}, \\
    g^{x_2} &\equiv p_1^{e_{21}} p_2^{e_{22}} \cdots p_k^{e_{2k}} \pmod{p}, \\
    &\;\;\vdots \\
    g^{x_m} &\equiv p_1^{e_{m1}} p_2^{e_{m2}} \cdots p_k^{e_{mk}} \pmod{p}.
\end{align*}
Since \(g\) is a generator of \(\mathbb{Z}_p^\times\), each factor base prime \(p_i\) has a discrete logarithm \(r_i = \log_g(p_i) \pmod{p-1}\), meaning \(p_i \equiv g^{r_i} \pmod{p}\).
Taking discrete logarithms modulo \(p-1\) transforms the multiplicative relations into a system of linear congruences:
\begin{align*}
    x_1 &\equiv e_{11} r_1 + e_{12} r_2 + \dots + e_{1k} r_k \pmod{p-1}, \\
    x_2 &\equiv e_{21} r_1 + e_{22} r_2 + \dots + e_{2k} r_k \pmod{p-1}, \\
    &\;\;\vdots \\
    x_m &\equiv e_{m1} r_1 + e_{m2} r_2 + \dots + e_{mk} r_k \pmod{p-1}.
\end{align*}
Writing this linear system in matrix form as \(E \mathbf{r} \equiv \mathbf{x} \pmod{p-1}\), where \(E = (e_{ji}) \in \mathbb{Z}^{m \times k}\), \(\mathbf{r} = (r_1, \dots, r_k)^\top\), and \(\mathbf{x} = (x_1, \dots, x_m)^\top\), we observe that the exponent matrix \(E\) is \emph{sparse}.
Indeed, each relation represents an integer \(g^{x_j} \pmod{p} < p\), and since any positive integer less than \(p\) can have at most \(\lfloor \log_2 p \rfloor\) prime factors (counted with multiplicity), each row of \(E\) has at most \(\log_2 p\) non-zero entries.
Since \(k = \pi(B)\) is subexponential in \(\ln p\), \(\log_2 p \ll k\), meaning that the vast majority of entries in every row are zero.

Because the matrix \(E\) is very sparse, we do not need to use dense Gaussian elimination, which would require \(O(k^3)\) operations.
Instead, specialized sparse linear algebra algorithms (such as the Wiedemann or Lanczos algorithms, or structured Gaussian elimination) can solve the system in \(O(k^2 \cdot \log p)\) operations, or roughly \(O(k^2)\) field operations, making the linear algebra phase significantly more efficient.
Once we collect \(m > k\) linearly independent relations, we solve for the factor base discrete logarithms \(r_1, \dots, r_k\).

\paragraph{Stage 2: Computing the Target Discrete Logarithm}
To find \(x = \log_g(h)\), we select random exponents \(s \in \mathbb{Z}_{p-1}\) and test whether
\[
    h \cdot g^s \pmod{p}
\]
is \(B\)-smooth.
Once such an exponent \(s\) is found, the residue factors over the factor base:
\[
    h \cdot g^s \equiv p_1^{f_1} p_2^{f_2} \cdots p_k^{f_k} \pmod{p}.
\]
Taking discrete logarithms with base \(g\) on both sides:
\[
    \log_g(h) + s \equiv f_1 r_1 + f_2 r_2 + \dots + f_k r_k \pmod{p-1}.
\]
Substituting the precomputed logarithms \(r_1, \dots, r_k\) yields the target discrete logarithm:
\[
    x = \log_g(h) \equiv \left(\sum_{i=1}^k f_i r_i\right) - s \pmod{p-1}.
\]

\subsection*{Smoothness Probability and Complexity Analysis}

To analyze the complexity of index calculus, we introduce the subexponential \(L\)-notation.

\begin{definition}[Subexponential \(L\)-Notation]
    For parameters \(0 \leq \alpha \leq 1\) and \(c > 0\), define
    \[
        L_n[\alpha, c] = \exp\left( (c + o(1)) (\ln n)^\alpha (\ln \ln n)^{1 - \alpha} \right).
    \]
\end{definition}

When \(\alpha = 0\), \(L_n[0, c] = (\ln n)^{c + o(1)}\) is polynomial in \(\ln n\).
When \(\alpha = 1\), \(L_n[1, c] = n^{c + o(1)}\) is exponential in \(\ln n\).
For \(0 < \alpha < 1\), \(L_n[\alpha, c]\) represents a subexponential growth rate.

The proportion of \(B\)-smooth integers up to \(p\) is given by the Canfield--Erd\H{o}s--Pomerance theorem.

\begin{theorem}[Canfield--Erd\H{o}s--Pomerance~\cite{CanfieldErdosPomerance1983}]
    Let \(u = \frac{\ln p}{\ln B}\).
    Then a uniformly chosen integer in \(\{1, \dots, p\}\) is \(B\)-smooth with probability
    \[
        u^{-u + o(u)} = \exp\left( (-1 + o(1)) \frac{\ln p}{\ln B} \ln\left( \frac{\ln p}{\ln B} \right) \right).
    \]
\end{theorem}

If we choose the smoothness bound \(B\) as
\[
    B = L_p\left[\frac{1}{2}, c_1\right] = \exp\left( (c_1 + o(1)) \sqrt{\ln p \ln \ln p} \right),
\]
then
\[
    u = \frac{\ln p}{\ln B} = \left(\frac{1}{c_1} + o(1)\right) \sqrt{\frac{\ln p}{\ln \ln p}},
\]
and
\[
    \ln u = \left(\frac{1}{2} + o(1)\right) \ln \ln p.
\]
Substituting \(u\) and \(\ln u\) into the smoothness probability:
\[
    u^{-u} = \exp(-u \ln u) = \exp\left( -\left(\frac{1}{2c_1} + o(1)\right) \sqrt{\ln p \ln \ln p} \right) = L_p\left[\frac{1}{2}, -\frac{1}{2c_1}\right].
\]
Hence, the expected number of trials to obtain a single \(B\)-smooth residue is \(L_p[1/2, 1/(2c_1)]\).

The size of the factor base is \(k = \pi(B) \leq B = L_p[1/2, c_1]\).
To collect \(k\) relations, the total relation-collection time is
\[
    k \cdot L_p\left[\frac{1}{2}, \frac{1}{2c_1}\right] = L_p\left[\frac{1}{2}, c_1 + \frac{1}{2c_1}\right].
\]
The linear algebra phase on a \(k \times k\) system requires \(O(k^2)\) operations (using sparse methods), which is \(L_p[1/2, 2c_1]\).
Balancing the relation collection and linear algebra stages by minimizing \(\max\{c_1 + \frac{1}{2c_1}, 2c_1\}\) gives
\[
    c_1 + \frac{1}{2c_1} = 2c_1 \implies c_1 = \frac{1}{\sqrt{2}}.
\]
This yields an overall heuristic running time of
\[
    L_p\left[\frac{1}{2}, \sqrt{2}\right] \quad \text{operations.}
\]

\begin{remark}[Cryptographic Consequences and the Number Field Sieve]
    While the basic index calculus algorithm runs in subexponential time \(L_p[1/2, \sqrt{2}]\), there exists an even more efficient variant known as the \emph{general number field sieve} (GNFS) for discrete logarithms.
    The number field sieve reduces the asymptotic complexity to
    \[
        L_p\left[\frac{1}{3}, \left(\frac{64}{9}\right)^{1/3}\right] \approx L_p\left[\frac{1}{3}, 1.923\right] \quad \text{operations.}
    \]
    Although the number field sieve is asymptotically faster and represents the state of the art for computing discrete logarithms over prime fields, its details rely on algebraic number theory and will not be covered in this course.
    Nonetheless, the existence of such \(L_p[1/3, c]\) subexponential attacks is precisely why key sizes for finite-field Diffie--Hellman must be substantially larger (e.g., \(2048\)--\(3072\) bits) to provide the same level of security as a \(256\)-bit elliptic curve group, where only generic algorithms (\(O(\sqrt{n})\)) apply.
\end{remark}
